\begin{table}[H]
  \centering
  \caption{任务对应}
  \label{tab:delivery-mapping}
  \begin{tabularx}{\linewidth}{lY}
    \toprule
    要点 & 报告和仓库中的证据 \\
    \midrule
    理论基础 & “基础”章说明量子态、酉门、SWAP、硬件耦合和噪声背景。\\
    算法设计 & “方法”章给出 NPQR 流程、动作评分、束搜索、局部修复和复放验证。\\
    核心代码 & “系统”章按功能结构说明算法层、服务层和证据层。\\
    实验结果 & “实验”章给出 10/20Q、选定 30/50Q 和分阶段耗时。\\
    工程实现 & “系统”和“部署”章说明 GitHub Pages、REST、MCP、命令行和服务器部署。\\
    AI 使用 & “AI 协作”章说明科研整理、前端、服务接口、测试部署和报告写作中的使用方式。\\
    过程管理 & 图 \ref{fig:project-timeline} 展示从基线、NPQR 迭代到公开交付的时间线。\\
    可复现性 & 附录列出安装、服务启动、示例编译和公开测试命令。\\
    \bottomrule
  \end{tabularx}
\end{table}
